\subsection{Norms from splitting fields and tame localization}\label{sec:second-two-tame}
In this subsection, we first show that, over a two-dimensional local field, the field norms from splitting fields generate $K_2(F)$ (Theorem~\ref{thm:second-two-surj}); in particular, the reduced norm on $K_2$ is surjective.
Under the hypotheses stated below, we then give localization formulas for $\SK_2(D)$ when $D/F$ is tamely ramified: in first residue characteristic zero (Proposition~\ref{prop:second-residue-zero}) and, for $\ell$-primary $D$ with $\ell$ different from the first residue characteristic, after $\ell$-localization (Proposition~\ref{prop:second-residue-zero-ell}).
Along the way we prove that $\SK_1(D)=0$ over two-dimensional local fields (Corollary~\ref{cor:review-two-local-sk1-zero}).
These results reduce the vanishing of $\SK_2$ over two-dimensional local fields to the wild case of Section~\ref{sec:second-two-wild}.

Let $F=F_h$ be an $h$-dimensional local field, with the fields $F_{h-1},\ldots,F_0$ as in the Conventions.
We write the characteristic sequence in the order $\operatorname{char} F_h,\ldots,\operatorname{char} F_1$.
For a central division algebra $D$ over $F$ and $\alpha=[D]\in\Br(F)$, let $N_2(\alpha)\subseteq K_2(F)$ be the norm subgroup of Definition~\ref{def:paper-transfer}.
We identify $K_2^M(F)$ with $K_2(F)$ by Matsumoto's theorem \cite{Matsumoto}.
The Morita-normalized splitting transfer $\theta_E$ satisfies $\Nrd_2\theta_E=N_{E/F}$, and hence $N_2(\alpha)\subset\Nrd_2K_2(D)$.

\begin{theorem}\label{thm:second-two-surj}
Let $F$ be a two-dimensional local field, let $D$ be a central division algebra over $F$, and put $\alpha=[D]$.
Then $N_2(\alpha)=K_2(F)$.
Consequently, $C_2(D/F)=0$.
\end{theorem}
\begin{proof}
We apply \cite[Proposition 5]{KatoExistence} to $\alpha\in H_m^2(F)=\Br(F)[m]$ with $m=\per(\alpha)$, using logarithmic coefficients for the characteristic-primary part.
Taking $n=q=r=2$ gives $q+r>n+1$, so the Milnor $K_2$-norms from finite extensions killing $\alpha$ generate $K_2(F)$.
\end{proof}

For the tame localization formulas we need the vanishing of $\SK_1$ over two-dimensional local fields, which we deduce from the following factorization of norm-one elements.

\begin{lemma}\label{lem:second-sl1-two-generators}
Let $F$ be a two-dimensional local field and $D$ a central division algebra over $F$.
Every element of $\operatorname{SL}_1(D)(F)$ is a product of two elements of $\operatorname{SL}_1(D)(F)$, each of which generates a separable maximal subfield of $D$.
\end{lemma}
\begin{proof}
The group $\operatorname{SL}_1(D)$ is smooth and geometrically connected, and open balls in \'etale coordinates for the nontrivial Henselian valuation on $F$ show that its $F$-points are Zariski dense.
The regular semisimple elements form a nonempty open subset $\Omega$, and an element of $\Omega(F)$ generates a separable maximal subfield.
For $g\in\operatorname{SL}_1(D)(F)$, the open set $\Omega\cap g\Omega^{-1}$ is nonempty, so by density it has an $F$-point $h$.
Then $g=h\cdot(h^{-1}g)$ with $h,h^{-1}g\in\Omega(F)$.
\end{proof}

\begin{corollary}
\label{cor:review-two-local-sk1-zero}
If $F$ is a two-dimensional local field and $D$ is a central division algebra over $F$, then $\SK_1(D)=0$.
\end{corollary}
\begin{proof}
Period equals index over a two-dimensional local field and over each of its finite extensions \cite[Corollary 5.1, Proposition 5.2]{Chipchakov2022}.
Moreover, Kato's duality and norm-image formula say that, for a Brauer class $\alpha$ of prime-power order, the subgroup $\alpha\cup F^\times\subset H^3(F,\Q/\Z(2))\simeq\Q/\Z$ has order $\per(\alpha)$, and the kernel of the cup product with $\alpha$ on $F^\times$ is the reduced-norm image.
Under the invariant map, restriction to a finite extension is multiplication by its degree.
The same formula with logarithmic coefficients applies to the characteristic-primary component \cite[\S3.1, Theorem 3(2)]{KatoLocalII}, \cite[Theorem 2]{KatoExistence}, \cite[Example 1.1(4)]{HuWu2020}.

The group $\SK_1(D)$ is killed by $\deg D$.
For the $\ell$-primary part, transfer makes restriction injective along an extension of degree prime to $\ell$ splitting the other Brauer components.
Morita equivalence then reduces to $\deg D=\ell^r$, and we induct on $r$ from the split case $r=0$.

By Lemma~\ref{lem:second-sl1-two-generators}, it is enough to treat a norm-one element $a$ that generates a separable maximal subfield.

Let $\widetilde L$ be the normal closure of $L=F(a)$ over $F$, let $G=\Gal(\widetilde L/F)$, and let $H$ be the subgroup fixing $L$.
We choose an $\ell$-Sylow subgroup $P$ of $G$ containing an $\ell$-Sylow subgroup of $H$, and pass to the extension $\widetilde L^P/F$ of degree prime to $\ell$.
The index and the degree of $L$ are unchanged, and restriction on $\SK_1$ remains injective.
A maximal subgroup of the $\ell$-group $P$ containing $H\cap P$ shows that, after this base change, $L$ contains a cyclic intermediate field $T/F$ of degree $\ell$.
We again denote the base field by $F$.

Put $B=Z_D(T)$ and $\alpha=[D]$.
Then $\deg B=\ell^{r-1}$, and period equals index, so $\per B=\ell^{r-1}$.
The element $b=\Nrd_{B/T}(a)$ satisfies $N_{T/F}b=1$; for a generator $\sigma$ of $\Gal(T/F)$, Hilbert 90 gives $b=c/\sigma(c)$ with $c\in T^\times$.
The order of $\alpha_T\cup c$ divides $\ell^{r-1}$, and the restriction of $\alpha\cup F^\times$ is also the subgroup of $\Q/\Z$ of order $\ell^{r-1}$.
We may therefore choose $u\in F^\times$ with $\alpha_T\cup(cu)=0$, and the norm-image formula then gives $y\in B^\times$ with $\Nrd_{B/T}(y)=cu$.
We choose $v\in D^\times$ such that conjugation by $v$ induces $\sigma$ on $T$.
Then
\[
 \Nrd_{B/T}([y,v])=c/\sigma(c)=b.
\]
Thus $a[y,v]^{-1}$ is a norm-one element of $B$, and its class in $K_1(B)$ vanishes by induction.
Since $[y,v]$ also vanishes in $K_1(D)$, so does the class of $a$.
We descend from the prime-to-$\ell$ extension and use the preceding factorization into two elements to conclude.
\end{proof}

We also need two lemmas, which are used again in Section~\ref{sec:second-three-coker}: surjectivity of reduction on $K_2$ for Azumaya algebras with division residue algebra over complete discrete valuation rings, and the vanishing of the cup product of the Brauer class with the image of the reduced norm.

\begin{lemma}\label{lem:second-three-azumaya-reduction}
Let $R$ be a complete discrete valuation ring with uniformizer $\pi$, and let $\widetilde D$ be an Azumaya algebra over $R$ whose residue algebra $\widetilde D/\pi\widetilde D$ is a division algebra.
Then $K_2(\widetilde D)\to K_2(\widetilde D/\pi\widetilde D)$ is surjective.
\end{lemma}
\begin{proof}
Put $I=\pi\widetilde D$, and carry out the argument of Lemma~\ref{lem:app-deep-unit-norm} over $R$.
Since $\widetilde D$ is Azumaya, its reduced trace pairing is perfect.
For units $s_i$ lifting a basis of the residue division algebra, the differential of $(x_i)\mapsto\prod_i[s_i,x_i]$ surjects onto $\ker\operatorname{Trd}$ on the residue, so the map to the smooth group $\operatorname{SL}_1(\widetilde D)$ is smooth at the identity.
Hensel's lemma expresses every norm-one principal unit as a product $\prod_i[s_i,x_i]$ with $x_i\in1+I$.
The relative group $K_1(\widetilde D,I)$ is generated by units in $1+I$, and, by the relative elementary relation $[1+bx]=[1+xb]$ used in the proof of Lemma~\ref{lem:app-orbit-stable-image}, the class of each $[s_i,x_i]$ is zero in $K_1(\widetilde D,I)$.
If a unit represents an element of the kernel of $K_1(\widetilde D,I)\to K_1(\widetilde D)$, then it has reduced norm one, so this map is injective.
The relative exact sequence $K_2(\widetilde D)\to K_2(\widetilde D/I)\to K_1(\widetilde D,I)\to K_1(\widetilde D)$ then makes the map $K_2(\widetilde D)\to K_2(\widetilde D/I)$ surjective.
\end{proof}

\begin{lemma}\label{lem:second-three-containment}
Let $K$ be a field, let $A$ be a central simple algebra over $K$, and put $\alpha=[A]$.
\begin{enumerate}
\item If $\operatorname{char}K=0$, then $\alpha\cup\Nrd_2(x)=0$ in $H^4(K,\Q/\Z(3))$ for every $x\in K_2(A)$.
If $\operatorname{char}K=p>0$, the same holds for the prime-to-$p$ components.
\item If $\operatorname{char}K=p>0$ and $A$ is a division algebra, then the $p$-primary component of $\alpha\cup\Nrd_2(x)$ vanishes for every $x\in K_2(A)$.
\end{enumerate}
\end{lemma}
\begin{proof}
(1) In characteristic zero, we use the \'etale motivic descent of Kahn--Levine \cite[\S6.9]{KahnLevine2010}.
The weight-two edge map of this descent is the reduced norm: this follows by passing to a splitting Severi--Brauer function field and using injectivity of restriction on $K_2(K)$.
By the module structure and the formula $d_2^A(1)=\alpha$ in \cite[Proposition 6.9.1]{KahnLevine2010}, the image of the edge map lies in the kernel of $d_2^A(x)=\alpha\cup x$.
In positive characteristic with invertible coefficients, the same calculation applies because passage to the perfect closure leaves the relevant Galois cohomology unchanged; here we also use the compatibility of $\Nrd_2$ with base change.

(2) Choose a complete Cohen discrete valuation ring $R$ with residue field $K$ and uniformizer $\pi$, and an Azumaya algebra $\widetilde D/R$ whose residue is $A$ \cite[Expos\'e XXIV, Proposition 8.1]{SGA3}.
By Lemma~\ref{lem:second-three-azumaya-reduction}, we can lift $x\in K_2(A)$ to $\widetilde x\in K_2(\widetilde D)$; let $y\in K_2(\operatorname{Frac}R)$ be its generic-point reduced norm.
Then $y$ comes from $K_2(R)$ and reduces to $\Nrd_2(x)$.
Indeed, after restriction to the discrete valuation ring of the generic point of the special fiber of the relative Severi--Brauer variety, where the algebra splits, the boundary of $y$ vanishes; since restriction on residue multiplicative groups is injective, the boundary of $y$ itself vanishes, and $y$ lifts to $K_2(R)$ by localization.
There the reduction on $K_2$ agrees with $s_\pi=\partial(\{\pi\}\cdot-)$ at the generic point, so two integral lifts with the same generic class have the same reduction; since $K_2(K)\to K_2(K(\SB(A)))$ is injective \cite[Corollaries 5.6--5.7]{Suslin1987}, the reduction of $y$ is $\Nrd_2(x)$.
The class $y$ has zero boundary, so it is a sum of unit symbols in $K_2(\operatorname{Frac}R)$.
Indeed, by Matsumoto's theorem it suffices to treat the symbols $\{\pi,1-\pi c\}$ with $c\in R$.
If $c\in R^\times$, the Steinberg relation gives $\{\pi,1-\pi c\}=-\{c,1-\pi c\}$.
If $c\in\pi R$, we write $1-\pi c=(1-\pi(1+c-\pi c))/(1-\pi)$ and apply the unit case to both factors.
Kato's unramified map $H^r(K)\hookrightarrow H^r(\operatorname{Frac}R)$ sends Artin--Schreier--Witt characters to unramified characters and sends products with unit symbols to products with arbitrary unit lifts \cite[\S5.1, pp.~55--56]{Kurihara2000}.
Thus $\alpha\cup\Nrd_2(x)$ maps to $[\widetilde D_{\operatorname{Frac}R}]\cup y$.
Part~(1), applied over the field $\operatorname{Frac}R$ of characteristic zero, makes this product zero.
\end{proof}

\begin{proposition}\label{prop:second-residue-zero}
Let $k$ be a local field of dimension $r\ge1$ with $\operatorname{char}k=0$, and let $F=k((t))$.
Let $B$ be the residue division algebra of a central division algebra $D/F$ (tamely ramified, since $\operatorname{char}k=0$), and let $L$ be the center of $B$.
The group $G=\Gal(L/k)$ is cyclic.
If $\SK_1(B)=0$, then
\[
 \SK_2(D)\simeq
 \ker\!\left(K_2(B)_G
   \xrightarrow{N_{L/k}\Nrd_{B/L,2}}K_2(k)\right).
\]
If $\Nrd_{B/L,2}$ is surjective, then $\SK_2(B)_G\to\SK_2(D)$ is surjective.
\end{proposition}

We keep the notation of Proposition~\ref{prop:second-residue-zero}, whose proof follows the next two lemmas.
In the Witt decomposition $[D]=a+(\chi,t)$ with $a\in\Br(k)$ and $\chi\in H^1(k,\Q/\Z)$, we have $L=k(\chi)$ and $[B]=a_L$.
Let $R$ be the valuation ring of $D$.
By \cite[Theorem 1]{Zheglov2004} and its proof, the residue map $R\to B$ has a ring section $\iota\colon B\to R$, the coefficient embedding, which extends the fixed coefficient field $k\subset F=k((t))$; this justifies the identification $L\otimes_kF=L((t))$ used below.
Since $B$ is a separable $k$-algebra, the coefficient embedding over the fixed $k$ is obtained by lifting through the square-zero kernels of $R/J^{s+1}\to R/J^s$, where $J=J(R)$, and passing to the inverse limit.
Normalize $w$ on $D$ by $w(D^\times)=\Z$, so $w(t)=e$, where $e=[L:k]$.
Under the tame isomorphism $\Gamma_D/\Gamma_F\simeq\Gal(L/k)$, let $\sigma_L$ be the generator corresponding to the value $1$.
Since $[B]=a_L$ is restricted from $k$, the automorphism $\sigma_L$ lifts to a semilinear automorphism $\sigma$ of $B$.
Extending $\sigma$ to $B\otimes_kF$ and applying Skolem--Noether to the two resulting $F$-embeddings of $B\otimes_kF$ into $D$, we obtain $u\in D^\times$ with $u\,\iota(b)=\iota(\sigma(b))\,u$ for $b\in B$.
The action on the residue center gives $w(u)\equiv1\bmod e$, so after multiplying by a power of the central element $t$ we may assume that $u$ is a uniformizer of $D$.
With this uniformizer, $R=B[[u;\sigma]]$.
We write $\sigma_*$ for the automorphism of $K_*(B)$ induced by $\sigma$; since $\sigma^{[L:k]}$ is inner, $\sigma_*$ defines an action of $G$.

\begin{lemma}\label{lem:second-residue-zero-localization}
In this setup, let $i_*\colon K_*(B)\to K_*(R)$ be the map from closed support, and put $V_2=\ker(K_2(R)\to K_2(B))$.
Then $i_*=\iota_*\circ(1-(\sigma^{-1})_*)$, the group $V_2$ is uniquely divisible, and the localization sequence takes the form
\[
 0\longrightarrow K_2(B)_G\oplus V_2
 \xrightarrow{\ \iota_*+\mathrm{incl}\ } K_2(D)\xrightarrow{\partial_D}K_1(B)^G
 \longrightarrow0.
\]
\end{lemma}
\begin{proof}
The right $R$-module $B$ has the projective resolution $0\to uR\to R\to B\to0$ of length one.
The right linear isomorphism $R\to uR$, $r\mapsto ur$, twists the left $B$-action by $\sigma^{-1}$, since $\iota(b)ur=u\,\iota(\sigma^{-1}(b))r$.
Additivity therefore gives $i_*=\iota_*\circ(1-(\sigma^{-1})_*)$.
Since $\iota$ splits reduction, $K_2(R)=\iota_*K_2(B)\oplus V_2$, and $\iota_*$ is injective on $K_*(B)$.
Hence $\coker(i_*\colon K_2(B)\to K_2(R))=K_2(B)_G\oplus V_2$ and $\ker(i_*\colon K_1(B)\to K_1(R))=K_1(B)^G$.
The localization sequence $K_2(B)\to K_2(R)\to K_2(D)\to K_1(B)\to K_1(R)$ now gives the displayed sequence.
Finite-coefficient rigidity in residue characteristic zero makes $V_2$ uniquely divisible \cite[\S3, Corollary 2 to Theorem 2]{SuslinYufryakov}.
\end{proof}

\begin{lemma}\label{lem:second-residue-zero-norms}
In this setup, let $U\subset D$ be an unramified maximal subfield.
Then
\[
 \partial_F\circ\Nrd_{D,2}=\Nrd_{B/L,1}\circ\partial_D\quad\text{inside }\bar U^\times,
\]
and, for $x\in K_2(B)$,
\[
 \Nrd_{D,2}(\iota_*x)=\res_{F/k}N_{L/k}\Nrd_{B/L,2}(x).
\]
Moreover, $\res_{F/k}\colon K_2(k)\to K_2(F)$ is split injective.
\end{lemma}
\begin{proof}
Such a field $U$ exists: lift a separable maximal subfield $K$ of $B$ to $\iota(K)((t))$.
Let $\mathcal O_U$ be the valuation ring of $U$ and $\bar U$ its residue field, a maximal subfield of $B$ containing $L$.
Restriction of scalars from $R$ to $\mathcal O_U$ is compatible with the localization sequences, and after the split Morita equivalence it is $\res_{U/F}\Nrd_{D,2}$ at the generic point and $\res_{\bar U/L}\Nrd_{B/L,1}$ on closed support.
Since $e(U/F)=1$, the field-side boundary $K_2(U)\to K_1(\bar U)$ composed with $\res_{U/F}$ is $\partial_F$ followed by restriction to the residue field $\bar U$, and the resulting localization diagram gives the first identity.

For the second identity, put $M=L((t))$ and $C=Z_D(M)$.
By Skolem--Noether, $\sigma^{[L:k]}=\operatorname{Ad}(\beta)$ on $B$ for some $\beta\in B^\times$, so $\beta^{-1}u^{[L:k]}$ commutes with $B$ and the centralizer of $L$ in $R=B[[u;\sigma]]$ is $B[[\beta^{-1}u^{[L:k]}]]$.
Hence $C=B\otimes_LM$: it contains this algebra, and the double centralizer theorem shows that the degrees agree.
We use the Morita $(C,D_M)$-bimodule $D$.
After this Morita equivalence, restriction of scalars from $D_M$ to $D$ is extension of scalars along $C\to D$, so the composite of extension of scalars along $C\to D$ with $\Nrd_{D,2}$ is $N_{M/F}\Nrd_{C/M,2}$ \cite[Proposition 6.6.1(1)]{KahnLevine2010}.
Therefore
\[
 \begin{aligned}
 \Nrd_{D,2}(\iota_*x)
 &=N_{M/F}\Nrd_{C/M,2}(x)\\
 &=N_{M/F}\res_{M/L}\Nrd_{B/L,2}(x)\\
 &=\res_{F/k}N_{L/k}\Nrd_{B/L,2}(x)
 \end{aligned}
\]
by this formula, compatibility of reduced norms with base change, and the base-change formula for $L\otimes_kF=M$.
Finally, specialization at $t$ splits $\res_{F/k}$.
\end{proof}

\begin{proof}[Proof of Proposition~\ref{prop:second-residue-zero}]
Let $\xi\in\SK_2(D)$.
By Lemma~\ref{lem:second-residue-zero-norms}, $\Nrd_{B/L,1}(\partial_D\xi)=\partial_F(1)=1$, and $\SK_1(B)=0$ gives $\partial_D\xi=0$.
By Lemma~\ref{lem:second-residue-zero-localization}, $\xi$ therefore lies in $K_2(B)_G\oplus V_2$.
Its $V_2$-component vanishes because $\deg D$ kills $\SK_2(D)$ and $V_2$ is uniquely divisible.
Thus $\SK_2(D)$ is the kernel of the reduced norm on $K_2(B)_G$, which Lemma~\ref{lem:second-residue-zero-norms} and the injectivity of $\res_{F/k}$ identify with the displayed kernel.

For the final assertion, suppose that $\Nrd_{B/L,2}$ is surjective; we first show that $\SK_1(B)=0$ holds automatically, so that the displayed identification is available.
Since $L/k$ is finite, $L$ is a local field of dimension $r$, and $\operatorname{char}L=0$.
If $r\le2$, then $\SK_1(B)=0$ by Theorem~\ref{thm:paper-local-norm} for $r=1$ and by Corollary~\ref{cor:review-two-local-sk1-zero} for $r=2$.
If $r\ge3$, then Lemma~\ref{lem:second-three-containment}(1) and the surjectivity of $\Nrd_{B/L,2}$ give $[B]\cup K_2(L)=0$.
Writing a Milnor symbol of length $r-1$ as the product of its first two entries with the rest, we obtain $[B]\cup K_{r-1}^M(L)=0$, and Kato's top-degree duality gives $[B]=0$ \cite[Theorem 2]{KatoExistence}; hence $B=L$ and again $\SK_1(B)=0$.
We now use cyclic $K_2$ Hilbert 90 \cite[Theorem 14.1]{MerkurjevSuslin1982}: the norm $N_{L/k}\colon K_2(L)_G\to K_2(k)$ is injective.
Hence the kernel equals $\ker(K_2(B)_G\to K_2(L)_G)$, and right exactness of coinvariants shows that this kernel is the image of $\SK_2(B)_G$.
\end{proof}

When the first residue characteristic is $p>0$, the same kernel formula holds after $\ell$-localization for $\ell\ne p$.
In the next proposition, $F$ is a complete discretely valued field with residue field $k$ of characteristic $p$, $D$ is an $\ell$-primary central division algebra over $F$ with residue division algebra $B$, and $L$ is the center of $B$.
Since $D/F$ is tamely ramified ($\ell\ne p$), the conjugation $z\mapsto u^{-1}zu$ by a uniformizer $u$ of $D$ restricts to a generator of $\Gal(L/k)$.
Another normalized choice $u'$ gives a unit $u'u^{-1}$ of the valuation ring $\mathcal O_D$ of $D$, whose residue acts by an inner automorphism, so the induced action $\bar\sigma_*$ of $\Gal(L/k)$ on $K_2(B)$ does not depend on the choice.
This is the action in the statement.

\begin{proposition}\label{prop:second-residue-zero-ell}
Let $F$ be a complete discretely valued field whose residue field $k$ is a local field of dimension at most two with $\operatorname{char}k=p>0$.
Let $\ell\ne p$ be a prime, let $D$ be an $\ell$-primary central division algebra over $F$ (so $D/F$ is tamely ramified), let $B$ be its residue division algebra, and let $L$ be the center of $B$.
\begin{enumerate}
\item If $\SK_1(B)=0$, then
\[
 \SK_2(D)_{(\ell)}\simeq
 \ker\!\left(K_2(B)_{\Gal(L/k)}
   \xrightarrow{N_{L/k}\Nrd_{B/L,2}}K_2(k)\right)_{(\ell)}.
\]
\item If $k$ is a nonarchimedean local field, then $\SK_2(D)_{(\ell)}=0$.
\end{enumerate}
\end{proposition}
The proof of Proposition~\ref{prop:second-residue-zero-ell} is given after Lemma~\ref{lem:second-residue-zero-ell-integral} below.

Let $\mathcal O_D$ be the valuation ring of $D$, so that $B=\mathcal O_D/J(\mathcal O_D)$.
By the structure of tame division algebras, the ramification index $e$ of $D/F$ equals $[L:k]$ and divides $\deg D$, so it is a power of $\ell$.
Let $u$ be a uniformizer of $D$, let $\pi$ be a uniformizer of $F$, and write $\sigma_*$ and $\bar\sigma_*$ for the actions of conjugation $z\mapsto u^{-1}zu$ on $K_2(\mathcal O_D)$ and on $K_2(B)$.
Let $r\colon K_2(\mathcal O_D)\to K_2(B)$ be the reduction and $i_*\colon K_2(B)\to K_2(\mathcal O_D)$ the map from closed support.
There is no coefficient field in general, so we lift through an unramified centralizer.
Let $M/F$ be the unramified lift of $L/k$, and let $\mathcal O_0$ be a maximal order in the centralizer of $M$ in $D$.
Then $\mathcal O_0$ is an Azumaya algebra over $\mathcal O_M$ with residue algebra $B$.

\begin{lemma}\label{lem:second-residue-zero-ell-integral}
In the situation of Proposition~\ref{prop:second-residue-zero-ell}, the reduction $r$ is surjective, $\ker r$ is uniquely $\ell$-divisible, and
\[
 \im i_*\cap\ker r=0.
\]
\end{lemma}
\begin{proof}
Lemma~\ref{lem:second-three-azumaya-reduction}, applied to the Azumaya algebra $\mathcal O_0$ over the complete discrete valuation ring $\mathcal O_M$, whose residue algebra is the division algebra $B$, shows that $K_2(\mathcal O_0)\to K_2(B)$ is surjective; hence so is $r$.

By the rigidity of Suslin and Yufryakov \cite[\S3, Theorem 2 and Corollary 2]{SuslinYufryakov}, the fiber
\[
 V=\fib\bigl(K(\mathcal O_D)\to K(B)\bigr)
\]
is uniquely $\ell$-divisible.
The image $C=\im\bigl(K_3(B)\to\pi_2V\bigr)$ is $\ell$-torsion-free.
Moreover, $C/\ell^a$ is a quotient of $K_3(B)/\ell^a$, which injects into $K_3(B;\Z/\ell^a)$.
At each residue step, rigidity and localization give $|K_i(B;\Z/\ell^a)|\le|K_i(\bar B;\Z/\ell^a)|\cdot|K_{i-1}(\bar B;\Z/\ell^a)|$ for the residue algebra $\bar B$.
Since $k$ has dimension at most two, starting from $K_3(B)$ takes at most two residue steps and leaves only the groups $K_i(\F_q;\Z/\ell^a)$ with $i\ge1$.
Their orders are bounded independently of $a$, and hence so is the order of $C/\ell^a$.
Since $C$ is $\ell$-torsion-free, $|C/\ell^a|=|C/\ell|^a$, so the uniform bound forces $C/\ell=0$, and $\ker r\simeq\pi_2V/C$ is uniquely $\ell$-divisible.

The localization sequence gives
\[
 i_*\circ r=1-\sigma_*,\qquad r\circ i_*=1-\bar\sigma_*,
\]
where the first identity follows from additivity for the sequence of bimodules $J(\mathcal O_D)\to\mathcal O_D\to B$, and the second comes from the derived tensor product with the residue algebra, whose $\operatorname{Tor}_0$ and $\operatorname{Tor}_1$ give the two terms.
If $y=i_*b$ lies in $\ker r$, then $\bar\sigma_*b=b$, and hence $\sigma_*y=y$.
Since $\pi u^{-e}\in\mathcal O_D^\times$ and $\pi$ is central, conjugation by $u^{-e}$ is conjugation by a unit of $\mathcal O_D$, so $\sigma^e$ is inner on $\mathcal O_D$ and $\sigma_*^e=1$.
Since $r$ is surjective, the first identity shows that $1+\sigma_*+\cdots+\sigma_*^{e-1}$ kills $\im i_*$; applied to the $\sigma_*$-invariant element $y$, this gives $e\,y=0$.
Since $e$ is a power of $\ell$ and $\ker r$ is uniquely $\ell$-divisible, $y=0$.
\end{proof}

\begin{proof}[Proof of Proposition~\ref{prop:second-residue-zero-ell}]
We compare boundaries by restriction of scalars to the unramified maximal subfield $U\subset D$ obtained by lifting a separable maximal subfield of $B$ through $\mathcal O_0$; the resulting diagram commutes on all of $K_2(D)$, which we also use below for arbitrary integral lifts.
After Morita equivalence, restriction of scalars to the valuation ring $\mathcal O_U$ of $U$ is $\res_{U/F}\Nrd_{D,2}$ at the generic point and $\res_{\bar U/L}\Nrd_{B/L,1}$ on closed support, where $\bar U$ is the residue field of $U$.
Since $e(U/F)=1$, the field-side boundary is restriction to the residue field, so the boundary of $\SK_2(D)$ lies in $\SK_1(B)$.
Let $s_\pi=\partial([\pi]\cdot-)\colon K_2(F)\to K_2(k)$ be the specialization.
On lifts from $\mathcal O_0$, the composite $K_2(\mathcal O_0)\to K_2(D)\xrightarrow{\Nrd_2}K_2(F)\xrightarrow{s_\pi}K_2(k)$ equals the reduction $K_2(\mathcal O_0)\to K_2(B)$ followed by $N_{L/k}\Nrd_{B/L,2}$.
This is computed through the centralizer of $M$, at the generic point of the special fiber of the relative Severi--Brauer scheme where the algebra splits, as in the proof of Lemma~\ref{lem:second-three-containment}(2), together with the compatibility of specialization and transfer along the unramified extension $M/F$ for the common uniformizer $\pi$.
Let $y\in\ker r$.
After restriction to $U$ and Morita equivalence, the reduction of $y$ is computed from the filtration of $\mathcal O_D/\pi\mathcal O_D$ by the $J^i$ $(0\le i<e)$, where $J=J(\mathcal O_D)$, whose graded pieces are twists of $B$; hence it is a sum of twists of $r(y)=0$, and $s_\pi(\res_U\Nrd_2y)=0$.
Taking the norm from $U$ gives $\deg D\cdot s_\pi(\Nrd_2y)=0$, and since $\ker r$ is $\ell$-divisible and $\deg D$ is a power of $\ell$, $s_\pi(\Nrd_2y)=0$.
As $K_2(\mathcal O_0)\to K_2(B)$ is surjective, every class of $K_2(\mathcal O_D)$ is a lift from $\mathcal O_0$ plus an element of $\ker r$, so the formula above holds on all of $K_2(\mathcal O_D)$ after $\ell$-localization.

(1) Suppose that $\SK_1(B)=0$.
By localization, an element of $\SK_2(D)$, whose boundary vanishes, lifts to $K_2(\mathcal O_D)$.
By Lemma~\ref{lem:second-residue-zero-ell-integral}, reduction gives an exact sequence
\[
 0\longrightarrow\ker r\longrightarrow\coker i_*\longrightarrow K_2(B)_{\Gal(L/k)}\longrightarrow0,
\]
and $\coker i_*$ embeds into $K_2(D)$ by localization.
Since $\ker r$ is uniquely $\ell$-divisible and $\deg D$ kills $\SK_2(D)$, the reduction is injective on $\SK_2(D)_{(\ell)}$.
Hence $\SK_2(D)_{(\ell)}$ lies in the $\ell$-localized kernel of $N_{L/k}\Nrd_{B/L,2}$ on $K_2(B)_{\Gal(L/k)}$.
Conversely, working in the $\ell$-localized groups, so that the lift and the correction below are taken $\ell$-locally, the reduced norm of a lift of an element of the residue norm kernel lies in
\[
 J'_F=\ker\partial_F\cap\ker s_\pi\subseteq K_2(F),
\]
and, since $s_\pi$ agrees with the reduction $r_F\colon K_2(\mathcal O_F)\to K_2(k)$ on $\ker\partial_F$ by the projection formula, $J'_F$ is the image of $\ker r_F$.
On the commutative side, $i_*r_F=0$ and $r_F$ is surjective, so $i_*=0$ and $K_2(\mathcal O_F)\to K_2(F)$ is injective; hence $J'_F\simeq\ker r_F$, which is uniquely $\ell$-divisible by the same rigidity and boundedness.
Scalar extension followed by $\Nrd_{D,2}$ is multiplication by $\deg D$ on $J'_F$, so we may divide this norm by $\deg D$ in $J'_F$, extend scalars, and subtract to obtain a lift with zero reduced norm.
This correction comes from a class in $K_2(\mathcal O_F)$ with zero reduction, so it does not change the residue class in $K_2(B)_{\Gal(L/k)}$.
This gives the reverse comparison after $\ell$-localization.

(2) If $k$ is a nonarchimedean local field, then so is $L$, and Theorem~\ref{thm:paper-local-norm} gives $\SK_1(B)=0$ and the isomorphism $\Nrd_{B/L,2}$.
Cyclic $K_2$ Hilbert 90 \cite[Theorem 14.1]{MerkurjevSuslin1982} makes $N_{L/k}$ injective on $K_2(L)_{\Gal(L/k)}$, so the kernel in (1) vanishes.
\end{proof}
